% Image credits for the photographs and AI illustrations of Book 5
% (University Biology, Year 3). Machine-readable license records live in
% images/book5/CREDITS.md; keep the two files in sync.
\chapter*{Créditos de las imágenes}

% \sloppy must be in force when the paragraph is BROKEN, hence the \par before
% \endgroup -- without it the group closes first and the tolerance is restored.
\begingroup\sloppy

Los dibujos de este libro son originales. Las fotografías se utilizan
bajo sus respectivas licencias libres, con agradecimiento a sus
autores:

\begin{itemize}
  \item Capítulo~\ref{ch:b3:genomics}: fotografía de Frederick Sanger
        --- dominio público (Wikimedia Commons).
  \item Capítulo~\ref{ch:b3:genetic-engineering}: arroz dorado junto a
        arroz común, International Rice Research Institute, 2011 ---
        CC~BY~2.0 (Wikimedia Commons).
  \item Capítulo~\ref{ch:b3:structural-biology}: fotografía de Max
        Perutz, Associated Press, 1962 --- dominio público (Wikimedia
        Commons).
  \item Capítulo~\ref{ch:b3:bacteriology}: retrato de Robert Koch,
        Wellcome Collection --- CC~BY~4.0 (Wikimedia Commons);
        fotografía de Alexander Fleming en su laboratorio, archivo del
        U.S. Navy Bureau of Medicine and Surgery --- sin restricciones
        de derechos conocidas (Wikimedia Commons).
  \item Capítulo~\ref{ch:b3:innate-immunity}: fotografía de Élie
        Metchnikoff, Bain News Service, Library of Congress --- dominio
        público (Wikimedia Commons).
  \item Capítulo~\ref{ch:b3:adaptive-immunity}: retrato grabado de
        Edward Jenner, 1807 --- dominio público (Wikimedia Commons).
  \item Capítulo~\ref{ch:b3:nervous-systems}: dibujo de una célula de
        Purkinje por Santiago Ramón y Cajal, 1899 --- dominio público
        (Wikimedia Commons).
  \item Capítulo~\ref{ch:b3:endocrinology}: fotografía de Frederick
        Banting y Charles Best, h.~1924 --- dominio público (Wikimedia
        Commons).
  \item Capítulo~\ref{ch:b3:molecular-evolution}: fotografía de Motoo
        Kimura, 1986, fotógrafo desconocido, vía S.~Kumar y
        T.~Gojobori, \emph{Molecular Biology and Evolution} (2023) ---
        CC~BY~4.0 (Wikimedia Commons).
  \item Capítulo~\ref{ch:b3:behavioural-ecology}: fotografía de Niko
        Tinbergen por Rob Mieremet, Anefo, 1973 --- CC0, Nationaal
        Archief (Wikimedia Commons).
  \item Capítulo~\ref{ch:b3:conservation-biology}: fotografía de
        Martha, la última paloma migratoria, Enno Meyer, 1912 ---
        dominio público (Wikimedia Commons); lobos en el parque
        nacional de Yellowstone, National Park Service, 1996 ---
        dominio público (Wikimedia Commons); Sirocco el kakapo,
        Department of Conservation, Nueva Zelanda, 2012 --- CC~BY~2.0
        (Wikimedia Commons).
\end{itemize}

Los enlaces a las fuentes completas y las direcciones de las licencias
de cada fotografía figuran en
\texttt{images/\allowbreak book5/\allowbreak CREDITS.md}, en el repositorio del libro.

\bigskip
Junto a las fotografías y los dibujos originales, unas noventa figuras
repartidas por los capítulos del volumen son ilustraciones generadas
con inteligencia artificial, producidas con la generación de imágenes
de OpenAI a partir de indicaciones escritas por los editores del
libro, y revisadas para comprobar su exactitud biológica antes de
incluirlas. La indicación completa de cada imagen generada consta en
\texttt{images/\allowbreak book5/\allowbreak ai/\allowbreak PROMPTS.md}, en el repositorio.
\par\endgroup
\clearpage
